% =============================================================================
%  Atto 7 — Bring-up & Validation  (~5 min, 3 slide)
%  The "save your evening" section: how we know the stack actually works.
% =============================================================================

\section{Bring-up \& validation}

% ---------------------------------------------------------------------------- %
\begin{frame}{The validation ladder}
  \begin{columns}[T,onlytextwidth]
    \column{0.5\textwidth}
      \textbf{Five rungs, each one isolating one ring.}
      \begin{enumerate}\setlength\itemsep{2pt}
        \item \textbf{Firmware standalone.}\\
              \footnotesize\texttt{cansend can0 0A\#…}\;forces RPM
              setpoints; \texttt{candump} watches \texttt{0x0B} reply.
        \item \textbf{ROS CAN stack.}\\
              \footnotesize Inject synthetic frames; check
              \texttt{report\_parser} emits typed msgs.
        \item \textbf{SystemInterface in mock.}\\
              \footnotesize Launch with
              \texttt{use\_mock\_hardware:=true} — controller works,
              hardware not yet involved.
        \item \textbf{Full SystemInterface, real hardware.}\\
              \footnotesize \texttt{ros2 control list\_controllers},\;
              \texttt{ros2 topic echo /odom}.
        \item \textbf{Nav2 goal $\to$ autonomous run.}\\
              \footnotesize End-to-end check, recordings on \texttt{rosbag2}.
      \end{enumerate}

    \column{0.5\textwidth}
      \centering
      \begin{tikzpicture}[node distance=2mm, every node/.style={font=\scriptsize}]
        \node[ll, minimum width=58mm] (a) {1. firmware in isolation};
        \node[ll, below=of a, minimum width=58mm] (b) {2. ROS CAN stack};
        \node[hl, below=of b, minimum width=58mm] (c) {3. SystemInterface (mock HW)};
        \node[hl, below=of c, minimum width=58mm] (d) {4. SystemInterface (real HW)};
        \node[hl, below=of d, minimum width=58mm] (e) {5. Nav2 autonomous run};
        \draw[arrll] (a)--(b);\draw[arrhl] (b)--(c);
        \draw[arrhl] (c)--(d);\draw[arrhl] (d)--(e);
      \end{tikzpicture}

      \vspace{0.6em}
      \footnotesize\textcolor{black!60}{Each rung deliberately exercises
      \emph{exactly one new ring} of the chain — when it breaks, you know
      where to look.}
  \end{columns}
\end{frame}

% ---------------------------------------------------------------------------- %
\begin{frame}[fragile]{Diagnostic toolbox}
  \begin{columns}[T,onlytextwidth]
    \column{0.5\textwidth}
      \textbf{Bus.}
      \begin{itemize}\setlength\itemsep{2pt}
        \item \texttt{candump -t a -ta can0}
        \item \texttt{cansniffer can0}
        \item scope CAN-H/CAN-L if nothing moves
      \end{itemize}

      \vspace{0.5em}
      \textbf{Firmware.}
      \begin{itemize}\setlength\itemsep{2pt}
        \item GPIO task toggles for timing
        \item ST-Link SWV for non-blocking traces
      \end{itemize}

    \column{0.5\textwidth}
      \textbf{ROS.}
      \begin{itemize}\setlength\itemsep{2pt}
        \item \texttt{ros2 control list\_controllers}
        \item \texttt{ros2 control list\_hardware\_interfaces}
        \item \texttt{ros2 topic hz /mivia\_rover/encoder\_rpms}
        \item \texttt{ros2 topic hz /odom}
        \item \texttt{rqt\_graph}
      \end{itemize}

      \vspace{0.6em}
      \footnotesize\textcolor{black!60}{Three views: bytes, tasks, graph.
      Together they localise most failures quickly.}
  \end{columns}
\end{frame}